% !TEX root = main.tex
% 在此填写开题报告的信息与正文；编译时使用 main.tex。
% 各命令中预填的是表格原有的提示文本，填写时请删除后再写入自己的内容。

% ==================== 基本信息 ====================
% 简况表和封面为固定布局，填写较长内容后请检查是否超出边框。
% 封面题目超出第一条横线时会自动续写到第二条横线；两行仍放不下时自动缩小题目字号。
\newcommand{\reporttitle}{}
\newcommand{\studentname}{}
\newcommand{\studentid}{}
\newcommand{\gender}{}
\newcommand{\supervisor}{}
\newcommand{\secondsupervisor}{（选填）}
\newcommand{\practicesupervisor}{}
\newcommand{\department}{软件学院}
\newcommand{\majorcategory}{电子信息 0854}
\newcommand{\majorfield}{软件工程 085405}
\newcommand{\reportdate}{}

% ==================== 勾选项 ====================
% 填写选项名称即可勾选，名称须与下方列出的完全一致，多个选项用英文逗号分隔；留空则不勾选。
% 培养方式：全日制、非全日制
\newcommand{\studymode}{}
% 专业学位类型：硕士、博士
\newcommand{\degreetype}{}
% 硕士成果形式：专题研究类论文、调研报告、案例分析报告、产品设计（作品创作）、方案设计、其他
\newcommand{\masterachievement}{}
\newcommand{\masterother}{}% 勾选“其他”时填写具体形式
% 博士成果形式：专题研究类论文、重大装备、仪器设备、其他硬件产品、软件产品、设计方案、技术标准、其他
\newcommand{\doctorachievement}{}
\newcommand{\doctorother}{}% 勾选“其他”时填写具体形式

\newcommand{\englishtitle}{}
\newcommand{\abstracttext}{}
% 中文主题词之间请用 \quad 分隔，如 {异常检测\quad 智能运维}；直接输入的空格会被忽略。
\newcommand{\keywords}{}
\newcommand{\englishkeywords}{}
\newcommand{\sourceoftopic}{%
  % 提示文本，填写时删除。
  工程博士必须填写研究所依托的重大工程项目，并标明项目编号。
}

% ==================== 报告正文 ====================
% 在各命令的花括号内填写内容，段落之间可用 \par。
% 第二至第七节的边框可随内容自动跨页。

% 二、选题依据
\newcommand{\topicbasis}{%
  % 提示文本，填写时删除。
  阐述该选题的研究意义，或工程设计的价值和意义，国内外概况和发展趋势，选题的先进性和实用性，技术难度及工作量。附上主要参考文献（列出作者、论文名称、期刊名称、出版年月）。
}

% 三、研究的创新之处
\newcommand{\innovation}{%
  % 提示文本，填写时删除。
  研究内容、拟采用的研究方法、技术路线等方面有哪些创新之处（博士必填，硕士选填）。
}

% 四、研究方法与技术路线
\newcommand{\methodroute}{%
  % 提示文本，填写时删除。
  拟采用的设计方法，技术路线，实践方案的可行性分析。
}

% 五、预期成果
\newcommand{\expectedresults}{%
  % 提示文本，填写时删除。
  预期达到的目标和取得的成果。
}

% 六、已取得的与学位申请成果研究内容相关的成果
\newcommand{\existingresults}{%
  % 提示文本，填写时删除。
  已发表或被接收发表的学术论文或其它相关创新性成果。
}

% 七、研究计划与进度安排
\newcommand{\researchplan}{%
  % 提示文本，填写时删除。
  应包括文献调研，工程设计，新工艺、新材料、新设备、新产品的研制和调试，实验操作，实验数据的分析处理，撰写论文等。
}
